\section{Introduction}

\textbf{We show that the standard instructed-lie benchmark identifies an instruction-induced behavioral contrast, not deception as a construct}---so its accuracy cannot establish that a detector responds to deception rather than to instruction-induced behavior. The setting is \emph{instructed roleplay}: because the instruction produces the deception, the signal's two candidate causes are not separately identifiable that way. \textbf{We do not claim that deception is undetectable}---we end by detecting it under a design that removes the confound (§\ref{sec:external_audit}). Nor do we estimate $\tau_D$, deception's own effect on the detector: it names only what a deception-attribution reading of the accuracy would need (§\ref{sec:identification}).

\noindent\fbox{\parbox{0.975\linewidth}{\small\textbf{Three questions, kept apart.}
Q1. Can a detector classify instructed-lie trials against truthful ones? \textbf{Yes}---we reproduce it and dispute none of it.
Q2. Does that signal arise from the deception rather than from following the instruction? \textbf{Unidentified}---the instruction moves both, so no accuracy measured this way separates them (§\ref{sec:identification}).
Q3. Will the detector fire when deception arises with no instruction to comply with, under incentive or goal conflict? \textbf{Not for the audited surface rule}, but prior work's probe battery \emph{does} separate deception on materials we built to make that test valid (§\ref{sec:external_audit}).
This paper is about Q2 and Q3: an \emph{identification} claim about a class of benchmarks, not a capability claim about detectors.}}

\begin{figure}[t]
\centering
\begin{tikzpicture}[
  >={Stealth[length=2mm]},
  var/.style={draw, rounded corners=2pt, inner sep=3.5pt, font=\small},
  every path/.style={->, thick},
]
% nodes
\node[var] (E) at (0,0) {$E$: deception instruction};
\node[var] (D) at (4.9,0.95) {$D$: deception};
\node[var] (C) at (4.9,-0.95) {$C$: compliance behavior};
\node[var] (S) at (9.9,0) {$S$: detector signal};
\node[var] (V) at (7.0,-2.5) {$V$: claim truth-value};
\node[var] (B) at (9.9,-2.5) {$B$: model belief};
% the unidentified pair, boxed, with the annotation in the box's empty centre
\draw[dashed, gray, rounded corners=3pt, -] (3.05,-1.62) rectangle (6.75,1.62);
\node[font=\scriptsize\itshape, gray, align=center] at (4.9,0)
  {$\mathrm{do}(E)$ moves both;\\not separately identifiable};
% the standard benchmark intervention, and the two that would identify the effect
\node[font=\scriptsize, align=center] at (1.75,1.18)
  {standard benchmark\\intervention: $\mathrm{do}(E)$};
\node[font=\scriptsize, align=left, anchor=west] at (-0.45,-2.06)
  {\textbf{Attribution requires}};
\node[draw, rounded corners=2pt, inner sep=1.6pt, font=\scriptsize, align=left, anchor=west]
  at (-0.45,-2.5) {\textbf{(A)} move $D$, hold $E$ $=$ \textbf{criterion 4}};
\node[font=\scriptsize, align=left, anchor=west] at (-0.45,-2.94)
  {\textbf{(B)} move $E$, hold $D$};
% edges
\draw (E) -- (D);
\draw (E) -- (C);
\draw (D) -- (S);
\draw (C) -- (S);
\draw (V) -- (S);
\draw (B) -- (S);
\end{tikzpicture}
\caption{\small\textbf{Why instructed-roleplay accuracy cannot attribute itself: the problem is identification, not classification.} The only variable such benchmarks manipulate is $E$, and no observable lies on one path but not the other. \textbf{What is needed is a different intervention, not better data} (Appendix Table~\ref{tab:roadmap}).}
\label{fig:dag}
\end{figure}

\paragraph{Contribution: three claims---\textsc{i.~identification}, \textsc{ii.~what a valid test requires}, \textsc{iii.~empirical diagnosis}---in descending order of how strongly they are warranted.} Appendix Table~\ref{tab:roadmap} maps each to its evidence; Figure~\ref{fig:ladder} shows how far up the evidence ladder they reach.

\noindent\textbf{Contribution~1 (\textsc{identification}).} \textbf{(1)~Framework.} Instructed benchmarks are informative about $\tau_E$, the \emph{instruction}'s effect, while deception-detection claims are about $\tau_D$---the \textbf{Elicitation--Construct Confounding Principle}, with behavioral deception detection as one instantiation (§\ref{sec:identification})---\emph{holds by structure}: \textbf{no experiment, ours or anyone's, could establish or refute this within an instructed design}.

\noindent\textbf{Contribution~2 (\textsc{what a valid test requires}).} \textbf{(2)~Protocol---an identification audit an auditor can apply to a benchmark they did not build}, and \emph{a deliverable, not a by-product}: five criteria separating instruction-following from surface behavior from knowledge conflict from \emph{recovery of its own annotation mechanism} passing for detection, one of them unsupplied by any instructed design (Figure~\ref{fig:ladder}; specified criterion by criterion in Appendix Table~\ref{tab:criteria}), and \textbf{decidable from a design's specification}---we decide them on ten designs in five groups, four outside deception (Appendix Table~\ref{tab:prior_work_criteria}), though a blinded panel re-deciding them from \emph{release} text alone reproduces only~(i) (EXP-AA). \textbf{What is new.} \textbf{The novelty is not the observation that confounding exists; it is turning a previously implicit causal-validity requirement into an auditor-executable protocol and showing that existing behavioral deception benchmarks fail it in empirically distinct ways.}

\noindent\textbf{Contribution~3 (\textsc{empirical diagnosis}).} \textbf{(3)~Demonstration: the confound is not hypothetical, it is active.} Established across eight models in four families, two claim sets and three detector paradigms, and an audit of both public rollout releases finds that \textbf{no audited public release supplies all five}---the closest cells failing in three distinct ways: an instrument-specific null, \textbf{annotation-channel leakage} ($+84.3$\,pp, the model echoing its own instruction), and \textbf{construct recovery} (§\ref{sec:experiments}, §\ref{sec:external_audit}; disclaimer~(2)). \textbf{(4)~A construct-valid benchmark: the critique does not imply undetectability.} So we built materials that do---and there, \textbf{prior work's own battery separates operationally verified deception at byte-identical elicitation} (pre-registered; \textbf{standing 1/5}; §\ref{sec:external_audit}), reaching \textbf{rung~4 of Figure~\ref{fig:ladder}, not rung~5}. The constraint is elementary but \emph{unenforced}---no benchmark we audit states requirement~(i)---so \textbf{the protocol and the audit, not the algebra, are what is new}.


\noindent \textbf{The full mapping---question $\to$ evidence $\to$ result $\to$ what it licenses, row by row, with pre-registration status marked---is Table~\ref{tab:roadmap}, in the appendix's reader's guide}, in three groups: \textbf{A.~Structural} (no experiment could establish or refute it), \textbf{B.~Empirical} (the confound is active), \textbf{C.~Attribution} (tested by criterion~4 alone, never by an instructed design). Its last row records what the table does \emph{not} establish: the \emph{absence} of a deception-specific signal.

\paragraph{Relation to prior work.} White-box probes read deception from activations~\citep{zou2023representation, burns2022discovering, azaria2023internal}; the black-box line's closest instance---our primary audit target---is \citet{pacchiardi2024catch}, and we audit both paradigms (§\ref{sec:td_vs_fd}). Three concurrent evaluations~\citep{kretschmar2025liarsbench, cooney2026didyoulie, hopkins2026liedetectors} find lie detectors fail to \textbf{generalize} beyond the regime they were fit in. That is a question about \textbf{predictive validity}; ours is the prior one, about \textbf{attributional validity}: instructed-lie accuracy is \textbf{not merely poorly generalizing but uninterpretable as evidence about deception}, because the intervention moves deception and compliance together---so generalization can be perfect while identification still fails. \textbf{Fixed elicitation is also not what we claim to introduce}---\citet{cooney2026didyoulie}'s trained organisms have it and go \emph{beyond} it, with $D$ \emph{set} rather than observed. \textbf{Our novelty is using it as the missing construct-validity criterion for reading instructed-benchmark accuracy}: we formalize the identification failure, turn it into an auditable protocol, and \textbf{show empirically that the confound accounts for substantial benchmark accuracy across three mechanistically distinct detector paradigms}---while auditing whether public artifacts supply the criterion at all. Appendix~\ref{app:related_work} places all five kinds against requirements~(i)--(v) (Appendix Table~\ref{tab:prior_work_criteria}).

\begin{center}\small
\setlength{\tabcolsep}{3pt}
\begin{tabular}{@{}p{0.11\linewidth}p{0.40\linewidth}p{0.42\linewidth}@{}}
\toprule
 & \textbf{Generalization} (concurrent work) & \textbf{Identification} (this paper) \\
\midrule
Validity & \textbf{Predictive} & \textbf{Attributional} \\
Asks & Does a detector fit in regime~A still work in regime~B? & Does accuracy in regime~A measure deception \emph{at all}? \\
Fails when & Accuracy drops out of distribution & One intervention moves two candidate causes \\
Fixed by & More regimes, more data, better detectors & \textbf{A different intervention} \\
\bottomrule
\end{tabular}
\end{center}

We begin with prior work's detector rather than our own, so no artifact of our reimplementation can carry the argument. \textbf{Four objects carry the argument}: Figure~\ref{fig:dag} shows why instructed accuracy cannot attribute itself, Figure~\ref{fig:ladder} how far the evidence reaches, \textbf{Table~\ref{tab:claim_ledger} every claim with its confirmatory experiment, primary metric, pre-registration and replication status}, and Table~\ref{tab:external_audit} what the audit found in every setting.

\noindent\fbox{\parbox{0.975\linewidth}{\small\textbf{The design that would settle the attribution question---stated in full, and then run.} A corpus that (i)~fixes the elicitation prompt, (ii)~grades $D$ independently of the condition label, (iii)~has \textbf{the same scenario independently realized under both behavioral outcomes}, (iv)~exposes only the channel a deployed detector could read, and (v)~grounds $D$ in evidence not definitionally equivalent to the detector's input---(i)--(v) operationalize criterion~4 (§\ref{sec:eval_controls}). We audit both public rollout releases---ten rollout sets and 35 further cells---and \textbf{none supplies all five}: (iii) is missing on the closest set of the first, and (v) on the second.
\textbf{This is falsifiable, and we say how}: a single public corpus supplying (i)--(v) in which a detector separates deception would overturn claim~(3). We ran that test on materials we built (§\ref{sec:external_audit}).}}

\noindent\textbf{What the accuracy aggregates}: \textbf{a heterogeneous mixture}---instruction-following (equalization takes 97.0\% to 43.3\%), surface correction markers (a parameter-free 14-pattern rule reaches 69--80\% with nothing fit), a component \emph{consistent with} knowledge conflict (70.0\% KNOWN vs.\ 57.6\% UNKNOWN, $p\!=\!0.0045$, stratified observational), and \emph{sometimes} deception-associated behavior (one of five targets, at byte-identical elicitation)---\textbf{and once the instruction confound is removed, what remains is neither stable across targets nor cleanly separable from surface behavior} (§\ref{sec:experiments}, §\ref{sec:external_audit}). Criteria~1--4 test these one at a time (Figure~\ref{fig:ladder}).

\paragraph{What we do not establish, in plain terms.} That deception is undetectable; that published detectors encode only instruction-following; that equalization shows a deception signal to be \emph{absent} rather than removed along with the confound; that the audited corpora contain no deception; or that the signal we do find at fixed elicitation is deception rather than a correlate of it. Each is stated precisely, with the evidence that bounds it, in §\ref{sec:eval_controls}.

\noindent\textbf{In one sentence}: \textbf{instructed-lie benchmarks answer Q1 and cannot answer Q2, because the instruction that creates the deception also creates the behavior a detector reads; answering Q2 requires materials in which deception varies while the elicitation is held byte-identical---and when we build them, prior work's own battery separates deception on one of five targets.}

